\section{Cross Attention 与 Self Attention}

\subsection{Cross Attention}

当 Query 向量和 Data 向量来自\textbf{两个不同的来源}时，这种注意力称为 \textbf{Cross Attention}（交叉注意力）。

\begin{figure}[H]
\centering
\includegraphics[width=0.95\textwidth]{slides-images/slide-030.jpg}
\caption{Cross Attention：Query 向量和 Data 向量来自不同的源。典型应用：机器翻译中解码器关注编码器、图像描述中文本生成器关注图像特征\protect\footnotemark}
\end{figure}
\footnotetext{Slides 第31页。}

Cross Attention 的典型应用场景：
\begin{itemize}
    \item \textbf{机器翻译}：解码器的状态（Query）关注编码器的隐藏状态（Key/Value）
    \item \textbf{图像描述}：生成的文本 token（Query）关注图像的 patch 特征（Key/Value）
    \item \textbf{多模态融合}：一种模态的表示（Query）从另一种模态中检索信息
\end{itemize}

\subsection{Self Attention}

当 Query、Key、Value \textbf{全部来自同一组输入向量}时，这种注意力称为 \textbf{Self Attention}（自注意力）。

\begin{figure}[H]
\centering
\includegraphics[width=0.95\textwidth]{slides-images/slide-032.jpg}
\caption{Self Attention：每个输入向量同时被投影为 Query、Key 和 Value。每个向量都能``看到''所有其他向量\protect\footnotemark}
\end{figure}
\footnotetext{Slides 第33页。}

\begin{importantbox}{Self Attention 的计算}
给定 $n$ 个输入向量 $X \in \mathbb{R}^{n \times d_{\text{in}}}$，self attention 通过三个可学习的权重矩阵将每个输入投影为 Query、Key 和 Value：

$$Q = X W_Q, \quad K = X W_K, \quad V = X W_V$$

然后执行标准的注意力计算：
$$Y = \text{softmax}\left(\frac{QK^T}{\sqrt{d_q}}\right) V$$

输出 $Y \in \mathbb{R}^{n \times d_v}$，与输入具有\textbf{相同数量}的向量。
\end{importantbox}

\subsection{本章小结}

Cross Attention 处理两个不同源的交互（如翻译中的源语言和目标语言），Self Attention 处理同一组向量内部的交互（如一个句子中词与词之间的关系）。Self Attention 是 Transformer 的核心构件。

%% ========== 置换等变性与位置编码 ==========

